\documentclass[11pt]{article}
\usepackage[margin=0.8in]{geometry}
\usepackage{amsmath,amssymb}
\usepackage[T1]{fontenc}
\usepackage{lmodern}
\usepackage[hidelinks]{hyperref}
\allowdisplaybreaks
\newcommand{\kw}[1]{\mathsf{#1}}
\newcommand{\step}[2]{\xrightarrow{#1,#2}}
\newcommand{\rulelabel}[1]{\quad\textsc{#1}}
\newcommand{\bg}{\operatorname{bg}}
\newcommand{\dom}{\operatorname{dom}}
\newcommand{\charge}{\operatorname{charge}}
\newcommand{\stable}{\operatorname{stable}}
\newcommand{\unstable}{\operatorname{unstable}}
\newcommand{\access}{\operatorname{access}}
\newcommand{\frm}[3]{(#1:#2,#3)}
\title{Cost-Aware Operational Semantics\\\large Consolidated Review Draft}
\author{}
\date{}
\begin{document}
\maketitle
\vspace{-2em}

\section{Scope and assumptions}
This draft consolidates the proposed operational rules for the language in
\emph{Leveraging Types and Typestate for Peripheral Cost Analysis}.
It retains a single processor, a stack of execution frames, and asynchronous
peripheral operations. It records both elapsed time and energy. It is a
source-level abstraction, rather than a cycle-accurate hardware model.

\paragraph{Timeout assumption.}
To retain the existing typing rule, timeout is an \emph{abstract exclusive
handoff}: once the fallback begins, the timed-out operation no longer accesses
the peripheral. The implementation of that handoff, including cancellation,
is deliberately outside this draft. An ordinary notification-only timeout
does not satisfy this assumption. Section~\ref{sec:timeout} isolates the rule
that depends on it.

\paragraph{Other assumptions.}
Expressions do not change peripherals or the clock. Function-local and
async-local variables are hygienically renamed; async code does not modify
the parent's variables. An async body accesses only its captured peripheral.
For this small model, async operations are not nested. Software operations are
atomic at the chosen source-language granularity; hardware state changes occur
between those operations. Hardware-only transitions have no separately charged
impulse energy; peripheral power accounts for their energy over elapsed time.
The hardware-only graph is acyclic, excluding infinite zero-time hardware runs.

\section{Configurations and notation}
\[
\Sigma=(M,P,C,S,K,t,E),
\qquad K\in\mathbb{Q}_{>0},\qquad t,E\in\mathbb{R}_{\geq0}.
\]
Here $M$ stores values and function definitions, $P$ records async operations,
$C$ is the execution stack (top on the left), $S$ maps peripherals to their
current states, $K$ is the clock frequency, $t$ is elapsed time, and $E$ is
accumulated energy. Peripheral declarations and hardware edges are stored in
a fixed declaration environment $\mathcal H$.
\begin{align*}
P(tid)&::=\kw{pending}(p,c)\mid\kw{running}(p)\mid\kw{done}(p),\\
F&::=\frm{tid}{p}{c}
\mid\frm{tid}{p}{[\,];c}
\mid\frm{tid}{p}{\kw{let}\ x=[\,]},\\
C&::=F::C\mid\epsilon.
\end{align*}
Commands additionally include runtime forms $\kw{skip}$, $\kw{endif}$, and
$\kw{ret}:=e$. A frame's peripheral field records its thread's captured
peripheral; the main thread may use a dummy value $\bot$.

Write $X=(M,P,C,S,K)$ for a configuration without counters. Local steps have
the form $X\step{q}{d}X'$, where $q$ is energy and $d$ is duration.
All cost and duration functions below are nonnegative. Computational steps
have positive duration; administrative steps have zero cost and duration.
The model supplies operation-specific functions $a_r$ and $d_r$.
For a simple CPU timing instance, $d_r(K)=N_r/K$; neither the rules nor the
energy model require energy to scale linearly with $K$.

\subsection{Background energy and ownership}
\[
\bg(S,K)=\sum_{p\in\dom(S)}\operatorname{power}(p,S(p),K),
\qquad
\charge(a,d,S,K)=a+d\,\bg(S,K).
\]
Define $\operatorname{owned}(P,p)$ to be the handle whose pending, running,
or done entry refers to $p$, if such a handle exists. There is at most one.
\[
\access(P,tid,p)
\iff
\bigl(\operatorname{owned}(P,p)\text{ is undefined}\bigr)
\lor\bigl(\operatorname{owned}(P,p)=tid\bigr).
\]
All software rules reading or changing $p$ require $\access(P,tid,p)$;
await rules are the exception because they operate on the handle.
Capture lasts until successful await or abstract timeout handoff.

\section{Expressions}
Expression evaluation returns a value and foreground energy. Its duration is
included in the enclosing command's duration $d_r(e,K)$.
\[
M,S,K\vdash e\Downarrow v\ ;\ a.
\]
\[
\frac{M(x)=v}{M,S,K\vdash x\Downarrow v\ ;\ 0}
\rulelabel{E-Var}
\qquad
\frac{}{M,S,K\vdash v\Downarrow v\ ;\ 0}
\rulelabel{E-Val}
\]
\[
\frac{
 M,S,K\vdash e_1\Downarrow v_1\ ;\ a_1
 \qquad M,S,K\vdash e_2\Downarrow v_2\ ;\ a_2
 \qquad v=v_1\mathbin{\mathrm{op}}v_2
}{
 M,S,K\vdash e_1\mathbin{\mathrm{op}}e_2\Downarrow v
 \ ;\ a_1+a_2+a_{\mathrm{op}}(K)
}\rulelabel{E-Binop}
\]
Set $a_{\mathrm{op}}=0$ to match the current expression typing rule, or add
the same operation term to that typing rule.

\section{Sequential commands and control flow}
\subsection{Assignment}
\[
\frac{
 M,S,K\vdash e\Downarrow v\ ;\ a_e
 \qquad d=d_{\mathrm{let}}(e,K)
 \qquad q=\charge(a_e+a_{\mathrm{let}}(K),d,S,K)
}{
 (M,P,\frm{tid}{p}{\kw{let}\ x=e}::C,S,K)
 \step{q}{d}
 (M[x\mapsto v],P,\frm{tid}{p}{\kw{skip}}::C,S,K)
}\rulelabel{E-Let}
\]

\subsection{Sequence}
The continuation stack removes the need for an unspecified append operation.
Both rules are administrative; all computational charges arise in $c_1,c_2$.
\[
\frac{}{
 (M,P,\frm{tid}{p}{c_1;c_2}::C,S,K)
 \step{0}{0}
 (M,P,\frm{tid}{p}{c_1}::\frm{tid}{p}{[\,];c_2}::C,S,K)
}\rulelabel{E-Seq}
\]
\[
\frac{}{
 (M,P,\frm{tid}{p}{\kw{skip}}::\frm{tid}{p}{[\,];c}::C,S,K)
 \step{0}{0}
 (M,P,\frm{tid}{p}{c}::C,S,K)
}\rulelabel{E-SeqDone}
\]

\subsection{Conditionals}
Let $c_{\kw{true}}=c_1$ and $c_{\kw{false}}=c_2$.
\[
\frac{
 M,S,K\vdash e\Downarrow b\ ;\ a_e
 \qquad d=d_{\mathrm{if}}(e,K)
 \qquad q=\charge(a_e+a_{\mathrm{if}}(K),d,S,K)
}{
 \begin{gathered}
 (M,P,\frm{tid}{p}{\kw{if}\ e\ \kw{then}\ c_1\ \kw{else}\ c_2}::C,S,K)
 \step{q}{d}\\
 (M,P,\frm{tid}{p}{c_b}::\frm{tid}{p}{\kw{endif}}::C,S,K)
 \end{gathered}
}\rulelabel{E-If}
\]
\[
\frac{}{
 (M,P,\frm{tid}{p}{\kw{skip}}::\frm{tid}{p}{\kw{endif}}::C,S,K)
 \step{0}{0}
 (M,P,\frm{tid}{p}{\kw{skip}}::C,S,K)
}\rulelabel{E-EndIf}
\]
The completed branch's $\kw{skip}$ must be removed together with the
$\kw{endif}$ marker; otherwise the branch would remain stuck on the stack.

\subsection{Peripheral matching}
Abbreviate $m=\kw{match}\ p\ \kw{with}\ st_1\to c_1|\cdots|st_n\to c_n$.
\[
\frac{
 S(p)=st_i\qquad\access(P,tid,p)
 \qquad d=d_{\mathrm{match}}(K)
 \qquad q=\charge(a_{\mathrm{match}}(K),d,S,K)
}{
 (M,P,\frm{tid}{r}{m}::C,S,K)
 \step{q}{d}
 (M,P,\frm{tid}{r}{c_i}::C,S,K)
}\rulelabel{E-Match}
\]

\section{Functions}
Store a function as $M(f)=(x_f,p_f,c_f,e_f)$. Let
\[
 (c^*,e^*)=\operatorname{instantiate}(c_f,e_f;
                         x_f\mapsto v,p_f\mapsto p).
\]
Instantiation substitutes the arguments capture-avoidingly and freshens
function-local bound variables. Functions have no externally visible memory
effects other than returning their value; peripheral effects are recorded in $S$.
\[
\frac{
 \begin{gathered}
 M(f)=(x_f,p_f,c_f,e_f)\qquad M,S,K\vdash e\Downarrow v\ ;\ a_e\\
 (c^*,e^*)=\operatorname{instantiate}(c_f,e_f;x_f\mapsto v,p_f\mapsto p)\\
 \access(P,tid,p)\qquad d=d_{\mathrm{call}}(e,K)
 \qquad q=\charge(a_e+a_{\mathrm{call}}(K),d,S,K)
 \end{gathered}
}{
 \begin{gathered}
 (M,P,\frm{tid}{r}{\kw{let}\ x=f(e,p)}::C,S,K)\step{q}{d}\\
 (M,P,\frm{tid}{r}{c^*;\kw{ret}:=e^*}
       ::\frm{tid}{r}{\kw{let}\ x=[\,]}::C,S,K)
 \end{gathered}
}\rulelabel{E-Call}
\]
\[
\frac{
 M,S,K\vdash e\Downarrow v\ ;\ a_e
 \qquad d=d_{\mathrm{ret}}(e,K)
 \qquad q=\charge(a_e+a_{\mathrm{ret}}(K),d,S,K)
}{
 \begin{gathered}
 (M,P,\frm{tid}{r}{\kw{ret}:=e}::\frm{tid}{r}{\kw{let}\ x=[\,]}::C,S,K)
 \step{q}{d}\\
 (M[x\mapsto v],P,\frm{tid}{r}{\kw{skip}}::C,S,K)
 \end{gathered}
}\rulelabel{E-Ret}
\]
The function cost summary must include these entry and return charges.
Alternatively, both may be zero-cost administrative operations, provided the
typing rules use exactly that convention.

\section{Peripheral and clock operations}
\subsection{Transition}
The energy $a_{\mathrm{tr}}$ includes the transitioning peripheral's energy
over the transition interval. Other peripherals continue drawing power.
\[
\frac{
 \begin{gathered}
 S(p)=s\qquad s\leq st\qquad\access(P,tid,p)\\
 d=d_{\mathrm{tr}}(p,s,st,K)\qquad
 q=a_{\mathrm{tr}}(p,s,st,K)+d\,\bg(S\setminus p,K)
 \end{gathered}
}{
 (M,P,\frm{tid}{r}{\kw{transition}(p,st)}::C,S,K)
 \step{q}{d}
 (M,P,\frm{tid}{r}{\kw{skip}}::C,S[p\mapsto st],K)
}\rulelabel{E-Transition}
\]
This rule handles stable and transient targets. Subsequent hardware steps
account for movement from a transient target through its hardware closure.

\subsection{Downgrade}
\[
\frac{
 \begin{gathered}
 S(p)=s\qquad(s\hookrightarrow st)\in D_p
 \qquad\access(P,tid,p)\\
 d=d_{\mathrm{tr}}(p,s,st,K)\qquad
 q=a_{\mathrm{tr}}(p,s,st,K)+d\,\bg(S\setminus p,K)
 \end{gathered}
}{
 (M,P,\frm{tid}{r}{\kw{downgrade}(p,st)}::C,S,K)
 \step{q}{d}
 (M,P,\frm{tid}{r}{\kw{skip}}::C,S[p\mapsto st],K)
}\rulelabel{E-Downgrade}
\]

\subsection{Clock update}
The new frequency takes effect at the end of the switch interval.
\[
\frac{
 K'>0\qquad d=d_{\mathrm{clock}}(K,K')
 \qquad q=a_{\mathrm{clock}}(K,K')+d\,\bg(S,K)
}{
 (M,P,\frm{tid}{p}{\kw{clock}\ K'}::C,S,K)
 \step{q}{d}
 (M,P,\frm{tid}{p}{\kw{skip}}::C,S,K')
}\rulelabel{E-Clock}
\]

\subsection{Hardware evolution}
\[
\frac{
 S(p)=s\qquad(s\xrightarrow{\mathrm{hw}}s')\in\mathcal H(p)
}{
 (M,P,C,S,K)\step{0}{0}(M,P,C,S[p\mapsto s'],K)
}\rulelabel{E-HW}
\]
Hardware steps can occur between software steps, independently of which
software thread is active. A hardware path consists of direct declared edges;
$\operatorname{HWReach}^{*}(A)$ denotes their reflexive transitive closure
from states in $A$. A transition's target is included in this closure.

\section{Asynchronous operations}
\subsection{Spawn}
Let $c^*=\operatorname{close}(c,M)$ snapshot the free ordinary variables of
$c$ by value and freshen its local bindings; peripheral identities remain
unchanged. The handle is bound in memory as well as in the typing context.
\[
\frac{
 \begin{gathered}
 \unstable(S(p))\qquad\operatorname{owned}(P,p)\text{ undefined}
 \qquad h\text{ fresh}\\
 c^*=\operatorname{close}(c,M)\qquad d=d_{\mathrm{spawn}}(K)
 \qquad q=\charge(a_{\mathrm{spawn}}(K),d,S,K)
 \end{gathered}
}{
 \begin{gathered}
 (M,P,\frm{tid}{r}{\kw{let}\ x=\kw{async}(p)\{c\}}::C,S,K)
 \step{q}{d}\\
 (M[x\mapsto h],P[h\mapsto\kw{pending}(p,c^*)],
       \frm{tid}{r}{\kw{skip}}::C,S,K)
 \end{gathered}
}\rulelabel{E-Async}
\]

\subsection{Scheduling and completion}
Scheduling pushes an async job above its suspended parent. Keep a running
entry in $P$ so that capture and handle lookup remain defined.
\[
\frac{P(h)=\kw{pending}(p,c)}{
 (M,P,C,S,K)\step{0}{0}
 (M,P[h\mapsto\kw{running}(p)],\frm{h}{p}{c}::C,S,K)
}\rulelabel{E-AsyncPreempt}
\]
Define $\operatorname{root}(h,C)$ to mean that $C$ is empty or its top frame
belongs to a thread other than $h$. This prevents completion from firing on
an intermediate $\kw{skip}$ inside a sequence or conditional.
\[
\frac{
 P(h)=\kw{running}(p)\qquad\operatorname{root}(h,C)
 \qquad\stable(S(p))
}{
 (M,P,\frm{h}{p}{\kw{skip}}::C,S,K)
 \step{0}{0}
 (M,P[h\mapsto\kw{done}(p)],C,S,K)
}\rulelabel{E-AsyncDone}
\]

\subsection{Await completion and waiting}
Write $w(h,n,c_{\mathrm{fb}})=
\kw{await}\ h\ \kw{timeout}[n]\{c_{\mathrm{fb}}\}$.
Source-level handle variables are resolved through $M$ to the corresponding
runtime handles $h$ when the await command becomes active.
\[
\frac{P(h)=\kw{done}(p)}{
 (M,P,\frm{tid}{r}{w(h,n,c_{\mathrm{fb}})}::C,S,K)
 \step{0}{0}
 (M,P\setminus h,\frm{tid}{r}{\kw{skip}}::C,S,K)
}\rulelabel{E-AwaitDone}
\]
\[
\frac{
 \begin{gathered}
 P(h)\in\{\kw{pending}(p,c),\kw{running}(p)\}
 \qquad n>0\qquad0<d\leq n\\
 q=d\bigl(\operatorname{power}_{\mathrm{wait}}(K)+\bg(S,K)\bigr)
 \end{gathered}
}{
 (M,P,\frm{tid}{r}{w(h,n,c_{\mathrm{fb}})}::C,S,K)
 \step{q}{d}
 (M,P,\frm{tid}{r}{w(h,n,c_{\mathrm{fb}})}::C,S,K)
}\rulelabel{E-Wait}
\]
The global step rule decrements the timer; the local waiting rule must not
decrement it a second time. Waiting uses the selected idle/sleep power model.
To exclude infinitely many shrinking wait intervals, choose a positive
timer quantum $\delta$ and set $d=\min(n,\delta)$ in this rule.

\subsection{Abstract timeout handoff}\label{sec:timeout}
Timeout has priority over further positive-duration execution once its timer
is zero. If completion has already been recorded, \textsc{E-AwaitDone} wins.
Under the exclusive-handoff assumption, use:
\[
\frac{
 \begin{gathered}
 P(h)\in\{\kw{pending}(p,c),\kw{running}(p)\}\\
 C=C_h\mathbin{+\!+}
   \bigl(\frm{tid}{r}{w(h,0,c_{\mathrm{fb}})}::C_0\bigr)
 \qquad\operatorname{allFramesOf}(C_h,h)
 \end{gathered}
}{
 (M,P,C,S,K)\step{0}{0}
 (M,P\setminus h,\frm{tid}{r}{c_{\mathrm{fb}}}::C_0,S,K)
}\rulelabel{E-AwaitTimeout}
\]
The prefix $C_h$ is empty for a pending operation; for a running operation,
it contains its execution frames. Removing it expresses the abstract handoff,
not a cancellation implementation. The actual current state $S(p)$ is retained;
timeout does not reset the peripheral to its state at spawn.

Timers are observed at software-step boundaries. Consequently, a step may
overshoot a deadline by its duration. A cost bound based on a timeout must
include this detection latency, or the model must assume steps are interruptible
at the deadline. This draft uses the boundary-observation convention.

\section{Global accounting and timers}
Define $\operatorname{drain}(C,d)$ structurally over all frames. For an await
command, including one stored in a suspended sequence continuation, replace
its residual timer by $\max(0,n-d)$ \emph{only if that await has started}.
To distinguish an unstarted source timeout from an active one, write the
runtime form as $w^{\bullet}(h,n,c_{\mathrm{fb}})$. Rules above using $w$ in
timer operations mean this runtime form.

Activation is administrative:
\[
\frac{M(x)=h}{
 \begin{gathered}
 (M,P,\frm{tid}{r}{\kw{await}\ x\ \kw{timeout}[n]\{c_{\mathrm{fb}}\}}::C,S,K)
 \step{0}{0}\\
 (M,P,\frm{tid}{r}{w^{\bullet}(h,n,c_{\mathrm{fb}})}::C,S,K)
 \end{gathered}
}\rulelabel{E-AwaitStart}
\]
\[
\operatorname{drain}\bigl(w^{\bullet}(h,n,c),d\bigr)
 =w^{\bullet}(h,\max(0,n-d),c).
\]
Unstarted source await commands and fallback bodies are not drained.
All active timers, including timers under a running async frame, are drained.
\[
\frac{
 (M,P,C,S,K)\step{q}{d}(M',P',C',S',K')
 \qquad C''=\operatorname{drain}(C',d)
}{
 (M,P,C,S,K,t,E)\longrightarrow
 (M',P',C'',S',K',t+d,E+q)
}\rulelabel{E-Step}
\]
The scheduler takes enabled administrative completion and timeout rules
before another positive-duration operation. It may otherwise choose a software
step, hardware step, or async dispatch as allowed above. A complete program
has a main $\kw{skip}$ frame, no remaining continuation, and an empty $P$.

\section{Declarations}
Declarations initialize the static environment and memory before execution;
their elaboration is not charged to the analyzed program.
\[
\frac{
 \mathcal H'=\mathcal H[p\mapsto(\mathcal S,s_0,\leq,D,H)]
 \qquad S'=S[p\mapsto s_0]
}{
 (\mathcal H,M,S,\kw{peripheral}\ p\{\mathcal S,s_0,\leq,D,H\})
 \longrightarrow(\mathcal H',M,S')
}\rulelabel{E-PeriphDecl}
\]
\[
\frac{M'=M[f\mapsto(x,p,c,e)]}{
 \begin{gathered}
 (\mathcal H,M,S,
 \kw{fn}\ f(x,p):(\tau_1,\vec s_1)\to(\tau_2,\vec s_2)
 \{c;\kw{ret}:=e\})\\
 \longrightarrow(\mathcal H,M',S)
 \end{gathered}
}\rulelabel{E-FunDecl}
\]

\section{Typing correspondence: items to check}
These are the interfaces the checking and symbolic systems must share with
the operational semantics. They are not additional operational rules.
\begin{enumerate}
\item Use identical foreground-energy and duration functions in both systems.
Interpret a background-cost term as energy over the operation's duration.
For state uncertainty $A_p$, a conservative background rate is
\[
 B(\Delta,K)=\sum_p\max_{s\in A_p}\operatorname{power}(p,s,K),
 \qquad S(p)\in A_p.
\]
If hardware may evolve, $A_p$ must cover its hardware closure. Mixed stable
and transient reachable sets must either be supported or explicitly excluded
by restrictions on peripheral declarations.
\item Keep structural sequence and branch-frame operations free. Function
summaries include any nonzero call and return charges. Spawn background
energy includes the captured peripheral unless that energy is already in
the spawn foreground cost.
\item The clock operation uses $a_{\mathrm{clock}}(K,K')$ and its duration.
Clock-dependent subcosts retain their own entry clocks. Replacing possible
clocks by their maximum requires monotone cost functions; otherwise keep
the alternatives, or require equal exit clocks at branch joins.
\item Async hardware and software effects preserve capture ownership. The
async body finishes only in a stable state. The fallback is checked from
\emph{every possible timeout state}, including hardware evolution and partial
async execution; using only the original spawn state is insufficient.
\item The existing await bound
$\max(Q_a,Q_{\mathrm{timeout}}+Q_{\mathrm{fb}})$ needs revision if it omits
waiting on success or partial work before timeout. A conservative accounting
partition is
\[
 Q_{\mathrm{await}}
 =Q_{a,\mathrm{remaining}}^{\mathrm{prefix}}
  +Q_{\mathrm{wait}}(n+\ell)+Q_{\mathrm{fb}},
\]
where $\ell$ bounds deadline-detection latency. Charge foreground async work
in the first term, all intervening background energy in the second, and
fallback energy in the third, without double charging. Work already accounted
for before await must be handled by the async resource invariant. This is a
conservative bound schema, not a proof that the current typing rule is sound.
\item The exclusive-handoff assumption is necessary for the existing fallback
ownership rule. If it is removed, the timeout typing rule must preserve capture.
\end{enumerate}

\paragraph{Target cost statement.}
For a well-typed initial configuration and a valid instantiated bound $Q$,
the intended theorem is
\[
 \Sigma_0\longrightarrow^*\Sigma_m
 \quad\Longrightarrow\quad E_m-E_0\leq Q
\]
for every finite execution prefix covered by the bound, together with
preservation of state membership and capture ownership. The corrected
operational rules alone do not establish this theorem; the typing obligations
above and an async accounting invariant must also be discharged.

\end{document}
